\chapter{Módulo 6 --- Visualización con ggplot2}

\begin{itemize}
\tightlist
\item
  \textbf{Duración estimada:} 1.5 h teoría + 2.5 h práctica
\item
  \textbf{Objetivos de aprendizaje.} Al finalizar, la persona podrá:

  \begin{enumerate}
  \def\labelenumi{\arabic{enumi}.}
  \tightlist
  \item
    \textbf{Explicar} la gramática de gráficos: datos, estéticas
    (\passthrough{\lstinline!aes!}), geometrías, escalas, facetas y
    temas.
  \item
    \textbf{Construir} gráficos de barras, líneas, dispersión,
    histogramas y cajas.
  \item
    \textbf{Seleccionar} el tipo de gráfico adecuado según la pregunta
    (comparar, tendencia, distribución, relación).
  \item
    \textbf{Producir} gráficos listos para reporte (títulos, etiquetas,
    formato de ejes) y exportarlos con
    \passthrough{\lstinline!ggsave()!}.
  \end{enumerate}
\end{itemize}

\section{Contenidos}

\begin{enumerate}
\def\labelenumi{\arabic{enumi}.}
\tightlist
\item
  La gramática de gráficos por capas.
\item
  \passthrough{\lstinline!aes()!}: \passthrough{\lstinline!x!},
  \passthrough{\lstinline!y!}, \passthrough{\lstinline!color!},
  \passthrough{\lstinline!fill!}, \passthrough{\lstinline!size!};
  estéticas fijas vs.~mapeadas.
\item
  Geometrías: \passthrough{\lstinline!geom\_col!},
  \passthrough{\lstinline!geom\_line!},
  \passthrough{\lstinline!geom\_point!},
  \passthrough{\lstinline!geom\_histogram!},
  \passthrough{\lstinline!geom\_boxplot!}.
\item
  Facetas: \passthrough{\lstinline!facet\_wrap()!}.
\item
  Escalas y formato: \passthrough{\lstinline!scales::label\_dollar()!},
  \passthrough{\lstinline!label\_percent()!}.
\item
  Etiquetas y temas: \passthrough{\lstinline!labs()!},
  \passthrough{\lstinline!theme\_minimal()!}.
\item
  Exportar: \passthrough{\lstinline!ggsave()!}.
\end{enumerate}

\section{Desarrollo}

\begin{lstlisting}[language=R]
library(tidyverse)
ventas  <- read_csv("datos/ventas.csv", show_col_types = FALSE) |>
  mutate(ingreso = precio * unidades)
tiendas <- read_csv("datos/tiendas.csv", show_col_types = FALSE)
ventas_det <- left_join(ventas, tiendas, by = "tienda_id")
\end{lstlisting}

\subsection{Comparar categorías: barras}

\begin{lstlisting}[language=R]
por_producto <- ventas |>
  summarise(ingreso = sum(ingreso, na.rm = TRUE), .by = producto)

ggplot(por_producto, aes(x = ingreso, y = fct_reorder(producto, ingreso))) +
  geom_col(fill = "steelblue") +
  scale_x_continuous(labels = scales::label_dollar(scale = 1e-6, suffix = " M")) +
  labs(title = "Ingreso anual por producto", x = "Ingreso (MXN)", y = NULL) +
  theme_minimal()
\end{lstlisting}

\begin{quote}
Barras \textbf{horizontales} y \textbf{ordenadas} facilitan leer nombres
largos y comparar magnitudes.
\end{quote}

\subsection{Tendencia: líneas}

\begin{lstlisting}[language=R]
semanal <- ventas_det |>
  summarise(ingreso = sum(ingreso, na.rm = TRUE), .by = c(fecha, canal))

g_tendencia <- ggplot(semanal, aes(fecha, ingreso, color = canal)) +
  geom_line(linewidth = 0.8) +
  scale_y_continuous(labels = scales::label_dollar()) +
  labs(title = "Ingreso semanal por canal, 2025", x = NULL, y = "Ingreso", color = "Canal") +
  theme_minimal()
g_tendencia
\end{lstlisting}

\subsection{Distribución: histograma y cajas}

\begin{lstlisting}[language=R]
ggplot(ventas, aes(unidades)) +
  geom_histogram(binwidth = 2, fill = "grey40", color = "white") +
  facet_wrap(~ producto, scales = "free_y") +
  labs(title = "Distribución de unidades semanales por producto")

ggplot(ventas_det, aes(region, unidades, fill = region)) +
  geom_boxplot(show.legend = FALSE) +
  labs(title = "Unidades por región", x = NULL)
\end{lstlisting}

\begin{quote}
Los \passthrough{\lstinline!NA!} de \passthrough{\lstinline!unidades!}
se descartan con una advertencia (``Removed rows\ldots{}''). Es
informativa, no un error.
\end{quote}

\subsection{Relación: dispersión}

\begin{lstlisting}[language=R]
por_tienda <- ventas_det |>
  summarise(unidades = sum(unidades, na.rm = TRUE),
            ingreso  = sum(ingreso,  na.rm = TRUE), .by = c(tienda_id, canal))

ggplot(por_tienda, aes(unidades, ingreso, color = canal)) +
  geom_point(size = 3) +
  geom_smooth(method = "lm", se = FALSE, color = "black", linewidth = 0.5) +
  labs(title = "Unidades vs. ingreso por tienda")
\end{lstlisting}

\subsection{Exportar}

\begin{lstlisting}[language=R]
ggsave("tendencia_canal.png", g_tendencia, width = 8, height = 4.5, dpi = 300)
\end{lstlisting}

\section{Actividades y ejercicios}

\begin{enumerate}
\def\labelenumi{\arabic{enumi}.}
\tightlist
\item
  Gráfico de barras apiladas al 100 \% de la participación de cada
  producto dentro de cada canal
  (\passthrough{\lstinline!position = "fill"!}).
\item
  Líneas de unidades mensuales por producto con
  \passthrough{\lstinline!facet\_wrap()!}.
\item
  Mejora un gráfico ``por defecto'' aplicando: orden, título que afirme
  la conclusión, ejes formateados y sin leyenda redundante.
\item
  Exporta tu mejor gráfico a PNG de 300 dpi.
\end{enumerate}

\subsection{Soluciones}

\begin{lstlisting}[language=R]
# 1
ventas_det |>
  summarise(ingreso = sum(ingreso, na.rm = TRUE), .by = c(canal, producto)) |>
  ggplot(aes(canal, ingreso, fill = producto)) +
  geom_col(position = "fill") +
  scale_y_continuous(labels = scales::label_percent()) +
  labs(title = "Mezcla de productos por canal", x = NULL, y = "Participación", fill = NULL)

# 2
ventas |>
  mutate(mes = floor_date(fecha, "month")) |>
  summarise(unidades = sum(unidades, na.rm = TRUE), .by = c(mes, producto)) |>
  ggplot(aes(mes, unidades)) +
  geom_line() +
  facet_wrap(~ producto, scales = "free_y") +
  labs(title = "Las ventas de todos los productos suben en noviembre y diciembre",
       x = NULL, y = "Unidades")
\end{lstlisting}

\section{Evaluación (formativa)}

Galería de pares: cada persona presenta un gráfico y recibe
retroalimentación con la lista de verificación: ¿responde una pregunta
clara?, ¿el tipo de gráfico es adecuado?, ¿ejes y unidades legibles?,
¿el título dice la conclusión?, ¿colores con propósito?

\section{Referencias verificadas}

\begin{itemize}
\tightlist
\item
  \textbf{Libro:} Wickham, H., Çetinkaya-Rundel, M., \& Grolemund, G.
  (2023). \emph{R for Data Science} (2.ª ed.). O'Reilly Media. ISBN
  978-1-4920-9740-2. Español: \url{https://davidrsch.github.io/r4ds-es/}
  {[}ES{]} {[}OA{]} --- caps. ``Visualización de datos'', ``Capas'',
  ``Análisis exploratorio'', ``Comunicación''.
\item
  \textbf{Web:} Wickham, H., Navarro, D., \& Pedersen, T. L.
  \emph{ggplot2: Elegant Graphics for Data Analysis} (3.ª ed., versión
  en línea). \url{https://ggplot2-book.org/} {[}EN{]} {[}OA en línea{]}
  \emph{{[}datos de edición impresa POR VERIFICAR{]}}
\item
  \textbf{Video:} freeCodeCamp.org. \emph{R Programming Tutorial --
  Learn the Basics of Statistical Computing} {[}Video{]}. YouTube.
  \url{https://youtu.be/_V8eKsto3Ug} {[}EN{]} --- secciones de gráficos
  (\passthrough{\lstinline!plot()!}, barras, histogramas, dispersión) en
  R base, útil para contrastar con ggplot2.
\end{itemize}
